\documentclass[10pt,a4paper]{amsart}
\usepackage{amsmath,amssymb,amsthm,graphicx}
\usepackage{verbatim}
\newenvironment{reference}{\comment}{\endcomment}
%\newenvironment{reference}{}{}
\newenvironment{slogan}{\comment}{\endcomment}
\newenvironment{history}{\comment}{\endcomment}

% For commutative diagrams we use Xy-pic
\usepackage[all]{xy}

% We use 2cell for 2-commutative diagrams.
\xyoption{2cell}
% Theorem environments.
%
\theoremstyle{plain}
\newtheorem{theorem}[subsection]{Theorem}
\newtheorem{proposition}[subsection]{Proposition}
\newtheorem{lemma}[subsection]{Lemma}

\theoremstyle{definition}
\newtheorem{definition}[subsection]{Definition}
\newtheorem{example}[subsection]{Example}
\newtheorem{exercise}[subsection]{Exercise}
\newtheorem{situation}[subsection]{Situation}

\theoremstyle{remark}
\newtheorem{remark}[subsection]{Remark}
\newtheorem{remarks}[subsection]{Remarks}

\numberwithin{equation}{subsection}

% Macros
%
\def\lim{\mathop{\mathrm{lim}}\nolimits}
\def\colim{\mathop{\mathrm{colim}}\nolimits}
\def\Spec{\mathop{\mathrm{Spec}}}
\def\Hom{\mathop{\mathrm{Hom}}\nolimits}
\def\Ext{\mathop{\mathrm{Ext}}\nolimits}
\def\SheafHom{\mathop{\mathcal{H}\!\mathit{om}}\nolimits}
\def\SheafExt{\mathop{\mathcal{E}\!\mathit{xt}}\nolimits}
\def\Sch{\mathit{Sch}}
\def\Mor{\mathop{\mathrm{Mor}}\nolimits}
\def\Ob{\mathop{\mathrm{Ob}}\nolimits}
\def\Sh{\mathop{\mathit{Sh}}\nolimits}
\def\NL{\mathop{N\!L}\nolimits}
\def\CH{\mathop{\mathrm{CH}}\nolimits}
\def\proetale{{pro\text{-}\acute{e}tale}}
\def\etale{{\acute{e}tale}}
\def\QCoh{\mathit{QCoh}}
\def\Ker{\mathop{\mathrm{Ker}}}
\def\Im{\mathop{\mathrm{Im}}}
\def\Coker{\mathop{\mathrm{Coker}}}
\def\Coim{\mathop{\mathrm{Coim}}}

% Boxtimes
%
\DeclareMathSymbol{\boxtimes}{\mathbin}{AMSa}{"02}

%
% Macros for moduli stacks/spaces
%
\def\QCohstack{\mathcal{QC}\!\mathit{oh}}
\def\Cohstack{\mathcal{C}\!\mathit{oh}}
\def\Spacesstack{\mathcal{S}\!\mathit{paces}}
\def\Quotfunctor{\mathrm{Quot}}
\def\Hilbfunctor{\mathrm{Hilb}}
\def\Curvesstack{\mathcal{C}\!\mathit{urves}}
\def\Polarizedstack{\mathcal{P}\!\mathit{olarized}}
\def\Complexesstack{\mathcal{C}\!\mathit{omplexes}}
% \Pic is the operator that assigns to X its picard group, usage \Pic(X)
% \Picardstack_{X/B} denotes the Picard stack of X over B
% \Picardfunctor_{X/B} denotes the Picard functor of X over B
\def\Pic{\mathop{\mathrm{Pic}}\nolimits}
\def\Picardstack{\mathcal{P}\!\mathit{ic}}
\def\Picardfunctor{\mathrm{Pic}}
\def\Deformationcategory{\mathcal{D}\!\mathit{ef}}
\usepackage{cleveref}
\crefname{originalreference}{original reference}{original references}
\Crefname{originalreference}{Original reference}{Original references}
\makeatletter
\newlabel{situation-shriek}{{0F42}{source}{}{sourcepointer-1}{}}
\newlabel{situation-shriek@cref}{{[originalreference][0][]0F42}{[1][0][]source}{}{}{}}
\makeatother
\crefname{originallemma}{original lemma}{original lemmas}
\Crefname{originallemma}{Original lemma}{Original lemmas}
\makeatletter
\newlabel{lemma-CM-shriek}{{0BV8}{source}{}{sourcepointer-2}{}}
\newlabel{lemma-CM-shriek@cref}{{[originallemma][0][]0BV8}{[1][0][]source}{}{}{}}
\makeatother
\crefname{originallemma}{original lemma}{original lemmas}
\Crefname{originallemma}{Original lemma}{Original lemmas}
\makeatletter
\newlabel{more-morphisms-lemma-flat-finite-presentation-CM-open}{{045U}{source}{}{sourcepointer-3}{}}
\newlabel{more-morphisms-lemma-flat-finite-presentation-CM-open@cref}{{[originallemma][0][]045U}{[1][0][]source}{}{}{}}
\makeatother
\crefname{originallemma}{original lemma}{original lemmas}
\Crefname{originallemma}{Original lemma}{Original lemmas}
\makeatletter
\newlabel{morphisms-lemma-openness-bounded-dimension-fibres}{{02FZ}{source}{}{sourcepointer-4}{}}
\newlabel{morphisms-lemma-openness-bounded-dimension-fibres@cref}{{[originallemma][0][]02FZ}{[1][0][]source}{}{}{}}
\makeatother
\crefname{originallemma}{original lemma}{original lemmas}
\Crefname{originallemma}{Original lemma}{Original lemmas}
\makeatletter
\newlabel{lemma-flat-shriek}{{0BV6}{source}{}{sourcepointer-5}{}}
\newlabel{lemma-flat-shriek@cref}{{[originallemma][0][]0BV6}{[1][0][]source}{}{}{}}
\makeatother
\crefname{originallemma}{original lemma}{original lemmas}
\Crefname{originallemma}{Original lemma}{Original lemmas}
\makeatletter
\newlabel{perfect-lemma-tor-independence-and-tor-amplitude}{{0C0V}{source}{}{sourcepointer-6}{}}
\newlabel{perfect-lemma-tor-independence-and-tor-amplitude@cref}{{[originallemma][0][]0C0V}{[1][0][]source}{}{}{}}
\makeatother
\crefname{originallemma}{original lemma}{original lemmas}
\Crefname{originallemma}{Original lemma}{Original lemmas}
\makeatletter
\newlabel{lemma-relative-dualizing-fibres}{{0BZZ}{source}{}{sourcepointer-7}{}}
\newlabel{lemma-relative-dualizing-fibres@cref}{{[originallemma][0][]0BZZ}{[1][0][]source}{}{}{}}
\makeatother
\crefname{originallemma}{original lemma}{original lemmas}
\Crefname{originallemma}{Original lemma}{Original lemmas}
\makeatletter
\newlabel{lemma-dualizing-module-CM-scheme}{{0AWT}{source}{}{sourcepointer-8}{}}
\newlabel{lemma-dualizing-module-CM-scheme@cref}{{[originallemma][0][]0AWT}{[1][0][]source}{}{}{}}
\makeatother
\title{For a flat morphism Cohen-Macaulayness at a point is equivalent to the shriek of the structure sheaf having one cohomology sheaf nearby}
\author{The Stacks Project Authors}
\date{stacks-project:duality.tex:local-snapshot-20261012; GNU FDL 1.2 or later}
\begin{document}
\maketitle
\noindent\textit{Editorial note.} All delivered mathematical text is the author's own, copied from the pinned Stacks Project chapter file duality.tex; no statement or proof step was rewritten, reordered or removed.\par\medskip
\noindent\textit{Editorial note.} Preamble provenance: the pinned chapter file opens with an input of the shared preamble and carries no preamble text of its own, so the project's single Stacks-wide preamble staged at sources/stacks-duality/preamble.tex is excerpted verbatim under the recipe fixed by LANE\_BRIEF 22.1, as three spans, lines 8-13 (verbatim package and the reference/slogan/history environments), lines 14-18 (Xy-pic setup) and lines 154-217 (the native theorem declarations and mathematical macros). Nothing was added, removed or redefined; the wrapper class and the amsmath/amssymb/amsthm/graphicx base packages are supplied by the assembler.\par\medskip
\noindent\textit{Editorial note.} The counted statement and proof contain no cross reference, no citation and no included file, so the delivered body is self contained.\par\medskip
\noindent\textit{Editorial note.} Canonical location of the counted proposition: Stacks tag 0C0Z (sources/stacks-duality/pinned-original-source.tex line 6284).\par\medskip
\noindent\textit{Allowed standard prerequisites (editorial).}\begin{itemize}
\item The Cohen-Macaulay criterion through the shriek of the structure sheaf. Same chapter, Lemma labelled lemma-CM-shriek, cited at proof lines 6295-6297 and 6304-6306. Gives one direction and the equality of the component dimensions.
\item Openness of the Cohen-Macaulay locus for a flat finite presentation morphism and boundedness of fibre dimensions. More on Morphisms, Lemma labelled more-morphisms-lemma-flat-finite-presentation-CM-open, cited at proof lines 6304-6306, and Morphisms, Lemma labelled morphisms-lemma-openness-bounded-dimension-fibres, cited at proof lines 6308-6310. Show the generic points are Cohen-Macaulay and bound the fibre dimensions.
\item Flatness of the single cohomology sheaf and the pullback along the fibre. Same chapter, Lemma labelled lemma-flat-shriek, cited at proof lines 6310-6312, and Derived Categories of Schemes, Lemma labelled perfect-lemma-tor-independence-and-tor-amplitude, cited at proof lines 6314-6316. Compute the pullback of the shriek to the fibre.
\item The dualizing complex of a fibre and the Cohen-Macaulay criterion for it. Same chapter, Lemmas labelled lemma-relative-dualizing-fibres and lemma-dualizing-module-CM-scheme, cited at proof lines 6316-6319. Identify the pullback as a dualizing complex and conclude.
\end{itemize}\medskip
\ifdefined\hypertarget\hypertarget{sourcepointer-1}{}\fi
\noindent\textit{Source reference.} Original reference 0F42: sources/stacks-duality/pinned-original-source.tex, source lines 3953--3953. Same-chapter pointer into the author environment situation at source line 3953; invoked at proof line None of the counted proof. It is retained only as a navigation pointer; every load-bearing step is inside the delivered span.\par\medskip
\ifdefined\hypertarget\hypertarget{sourcepointer-2}{}\fi
\noindent\textit{Source reference.} Original lemma 0BV8: sources/stacks-duality/pinned-original-source.tex, source lines 5978--5978. Same-chapter pointer into the author environment lemma at source line 5978; invoked at proof line 6296 of the counted proof. It is retained only as a navigation pointer; every load-bearing step is inside the delivered span.\par\medskip
\ifdefined\hypertarget\hypertarget{sourcepointer-3}{}\fi
\noindent\textit{Source reference.} Original lemma 045U: sources/stacks-chow/tags.txt, source lines 5449--5449. Cross-chapter standard prerequisite located in the project tag list staged in the local snapshot; invoked at proof line 6303 of the counted proof. It is retained only as a navigation pointer; every load-bearing step is inside the delivered span.\par\medskip
\ifdefined\hypertarget\hypertarget{sourcepointer-4}{}\fi
\noindent\textit{Source reference.} Original lemma 02FZ: sources/stacks-chow/tags.txt, source lines 3210--3210. Cross-chapter standard prerequisite located in the project tag list staged in the local snapshot; invoked at proof line 6307 of the counted proof. It is retained only as a navigation pointer; every load-bearing step is inside the delivered span.\par\medskip
\ifdefined\hypertarget\hypertarget{sourcepointer-5}{}\fi
\noindent\textit{Source reference.} Original lemma 0BV6: sources/stacks-duality/pinned-original-source.tex, source lines 5915--5915. Same-chapter pointer into the author environment lemma at source line 5915; invoked at proof line 6309 of the counted proof. It is retained only as a navigation pointer; every load-bearing step is inside the delivered span.\par\medskip
\ifdefined\hypertarget\hypertarget{sourcepointer-6}{}\fi
\noindent\textit{Source reference.} Original lemma 0C0V: sources/stacks-chow/tags.txt, source lines 15100--15100. Cross-chapter standard prerequisite located in the project tag list staged in the local snapshot; invoked at proof line 6313 of the counted proof. It is retained only as a navigation pointer; every load-bearing step is inside the delivered span.\par\medskip
\ifdefined\hypertarget\hypertarget{sourcepointer-7}{}\fi
\noindent\textit{Source reference.} Original lemma 0BZZ: sources/stacks-duality/pinned-original-source.tex, source lines 5130--5130. Same-chapter pointer into the author environment lemma at source line 5130; invoked at proof line 6315 of the counted proof. It is retained only as a navigation pointer; every load-bearing step is inside the delivered span.\par\medskip
\ifdefined\hypertarget\hypertarget{sourcepointer-8}{}\fi
\noindent\textit{Source reference.} Original lemma 0AWT: sources/stacks-duality/pinned-original-source.tex, source lines 6214--6214. Same-chapter pointer into the author environment lemma at source line 6214; invoked at proof line 6317 of the counted proof. It is retained only as a navigation pointer; every load-bearing step is inside the delivered span.\par\medskip
\begin{lemma}
\label{lemma-affine-flat-Noetherian-CM}
In Situation \ref{situation-shriek} let $f : X \to Y$ be a morphism
of $\textit{FTS}_S$. Let $x \in X$. If $f$ is flat, then
the following are equivalent
\begin{enumerate}
\item $f$ is Cohen-Macaulay at $x$,
\item $f^!\mathcal{O}_Y$ has a unique nonzero cohomology sheaf
in a neighbourhood of $x$.
\end{enumerate}
\end{lemma}

\begin{proof}
One direction of the lemma follows from Lemma \ref{lemma-CM-shriek}.
To prove the converse, we may assume $f^!\mathcal{O}_Y$ has a unique
nonzero cohomology sheaf. Let $y = f(x)$. Let $\xi_1, \ldots, \xi_n \in X_y$
be the generic points of the fibre $X_y$ specializing to $x$.
Let $d_1, \ldots, d_n$ be the dimensions of the corresponding
irreducible components of $X_y$. The morphism $f : X \to Y$ is Cohen-Macaulay
at $\eta_i$ by More on Morphisms, Lemma
\ref{more-morphisms-lemma-flat-finite-presentation-CM-open}.
Hence by Lemma \ref{lemma-CM-shriek} we see that
$d_1 = \ldots = d_n$. If $d$ denotes the common value, then $d = \dim_x(X_y)$.
After shrinking $X$ we may assume all fibres have dimension at most $d$
(Morphisms, Lemma \ref{morphisms-lemma-openness-bounded-dimension-fibres}).
Then the only nonzero cohomology sheaf $\omega = H^{-d}(f^!\mathcal{O}_Y)$
is flat over $Y$ by Lemma \ref{lemma-flat-shriek}.
Hence, if $h : X_y \to X$ denotes the canonical morphism, then
$Lh^*(f^!\mathcal{O}_Y) = Lh^*(\omega[d]) = (h^*\omega)[d]$
by Derived Categories of Schemes, Lemma
\ref{perfect-lemma-tor-independence-and-tor-amplitude}.
Thus $h^*\omega[d]$ is the dualizing complex of $X_y$ by
Lemma \ref{lemma-relative-dualizing-fibres}.
Hence $X_y$ is Cohen-Macaulay by
Lemma \ref{lemma-dualizing-module-CM-scheme}.
This proves $f$ is Cohen-Macaulay at $x$ as desired.
\end{proof}

\end{document}
